% ------------------------------------------------------------
% Page 75
% ------------------------------------------------------------

\begin{center}
    {\Large\bfseries Gatuzieres}
\end{center}

\vspace{0.45cm}

Après le Moulin d’Ayres nous traversons les hameaux de Salvin sac puis de la Bragouse et la
Bragousette pour arriver à Gatuzieres, dominé par les premiers lacets de la route du col de Perjuret.

\vspace{0.3cm}

Petit village où sont regroupées, vers les années 1900, une douzaine de familles avec sa mairie,
son temple construit après la Révolution vers 1840, et son église beaucoup plus ancienne
accompagné d’un très grand presbytère qui fut occupé par quelques sœurs « institutrices » de
l’école catholique Car il y avait, avant les lois de Jules Ferry, deux écoles, une protestante
avec un instituteur et l’autre catholique « tenue » par une ou deux sœurs. Elles laisseront la
place à l’école publique, construite au bas du village et maintenant mairie.

\vspace{0.3cm}

\textbf{Le « château » :} la plus belle propriété qui, de « château » n’en porte que le nom, encore qu’elle
possède un pigeonnier, en ruines aujourd’hui, signe de noblesse. On peut admirer une grange aux
dimensions imposantes avec ses six arcs en ogive dont la construction ne parait pas ancienne. A une
extrémité l’ancien logement du fermier ; celui du propriétaire à l’autre extrémité donne sur un
balcon couvert dont deux piliers proviennent, peut-être, des restes du château. \textbf{Le moulin du Mas
Pradès} --sans doute moulin du château-- fut en activité jusqu’aux années 1960; toutes les meules sont
encore en place. Tous deux étaient situés en aval du village.

\vspace{0.3cm}

A Gatuzieres vivaient quelques petits propriétaires disposant de maigres parcelles de terre et
des journaliers ou bergers engagés par les grosses fermes de la vallée et du Causse Méjean.
Jusqu’aux années 1960 chaque troupeau de brebis était conduit et gardé par un berger drapé
dans sa cape rigide et imperméable et accompagné de son grand parapluie bleu ou rouge,
visible de loin.

\vspace{0.3cm}

Il y avait aussi à Gatuzieres une recette postale avec sa cabine téléphonique. Mme. Dely
Martin tenait le bureau, installée derrière la banque et son grillage. Son mari distribuait le
courrier dans la commune . A sa mort c’est leur fille Alice qui remplaça son père. A 5h00 du
matin elle allait chercher lettres et plis apportés par « le courrier » l’autobus Meyrueis-Florac
qui s’arrêtait bien au dessus du village. Elle montait à Aures puis à Cabrillac, à pied bien sûr
passant par Plambel et Jontanels mais laissant de côté Malbosc desservi par Fraissinet de
Fourques !!! Il lui arrivait de rencontrer le facteur de Rousses qui, lui, avait en charge la
moitié de Cabrillac ( les maisons côté Nord ) !!!

\vspace{0.3cm}

La guerre de 1914 puis celle de 1939-1945 ainsi que la modernisation de l’agriculture
aggravèrent l’exode rural au point que dans les années 1970 il n’y eut plus qu’un seul habitant
permanent au village ; seuls, le moulin du Mas Pradès et le château « lou Castel », dont les
prés habitent désormais un camping, résistèrent.
Le village de Gatuzieres semblait condamné, tout comme Malbosc, Plambel, Jontanels et Cabrillac à
abriter des résidences secondaires préférables, somme toute, à des villages en ruines.

\vspace{0.3cm}

Mais en 1975-1980, un éleveur s’installa à Jontanels ; on vit, à Gatuzieres, se rouvrir
timidement et se réhabiliter quelques maisons. L’installation de la maison de retraite de
Meyrueis puis de petites entreprises ainsi que le Parc des Cévennes, accompagnés du
développement du tourisme avec ses campings et ses gîtes chez l’habitant et le retour de
quelques retraités redonnèrent vie à Gatuzieres où seule une maison en ruine dresse encore ses
pans de murs, sans toit, en plein village !

\vspace{0.35cm}

\begin{center}
    Mais en 2008 Gatuzieres est sauvé !!!
\end{center}
